% SOURCE_AUTHORITY: sga5_fr_workpass.tex, lines 1858--1988 (LF-normalized unit).
% SOURCE_SHA256: 791F4EFFC5E02832D5D77ED03518C8156D6F07E4C8238B03545DB93D883FBB28
\subsection*{1.1}
Denotamos por $S$ o espectro de um corpo. Salvo indicação em contrário, os $S$-esquemas considerados a seguir serão $S$-esquemas separados e de tipo finito.

Denotamos por $A$ um anel comutativo noetheriano, anulado por um inteiro inversível em $S$. Se $X$ é um $S$-esquema, denotamos por $D(X)=D(X,A)$ a categoria derivada dos complexos de $A$-módulos sobre $X$ (para a topologia étale), e por $D_{\ctf}(X)$ a subcategoria plena formada pelos complexos de dimensão Tor finita e cohomologia construtível. A categoria $D_{\ctf}(X)$ é uma subcategoria triangulada de $D(X)$. Se $E,F \in \ob D_{\ctf}(X)$, deduz-se facilmente das definições que $E \otimes^{L} F \in \ob D_{\ctf}(X)$. Segundo (SGA~$4^{1/2}$ Teorema de finitude 1.6, 1.7), tem-se também $\uRHom(E,F) \in \ob D_{\ctf}(X)$.

\subsection*{1.2}
Seja $f:X\to Y$ um $S$-morfismo. Tem-se trivialmente
\[
f^{*}D_{\ctf}(Y) \subset D_{\ctf}(X),
\]
e, segundo (loc. cit.),
\[
\R f_{*}D_{\ctf}(X) \subset D_{\ctf}(Y),\qquad
\R f_{!}D_{\ctf}(X) \subset D_{\ctf}(Y),\qquad
\R f^{!}D_{\ctf}(Y) \subset D_{\ctf}(X),
\]
(a inclusão relativa a $\R f_{!}$ utiliza apenas o teorema de finitude para os morfismos próprios (SGA~4~XIV) e o sorites (SGA~4~XVII~5.2.11); portanto, vale mais geralmente para todo morfismo compactificável cujas fibras tenham dimensão limitada superiormente). A seguir escreveremos às vezes $f_{*}$, $f_{!}$, $f^{!}$ em vez de $\R f_{*}$, $\R f_{!}$, $\R f^{!}$; as hipóteses de 1.1 implicam que esses funtores estão definidos sobre toda a categoria derivada.

\subsection*{1.3}
Para $f:X\to S$, poremos
\[
K_X=f^{!}A,
\]
e denotaremos por $D_X$ (ou $D$) o funtor $\uRHom(-,K_X)$. Segundo (SGA~$4^{1/2}$ Teorema de finitude 4.3), o complexo $K_X$ é dualizante: para $E\in\ob D_{\ctf}(X)$, tem-se canonicamente
\[
E \xrightarrow{\sim} DDE.
\]

\subsection*{1.4}
Sejam $f:X\to Y$ um $S$-morfismo, $L\in\ob D(X)$, $M\in\ob D(Y)$. Poremos
\[
\Hom_f^{(1)}(L,M)=\Hom(f^{*}M,L)\quad (\simeq \Hom(M,f_{*}L)),
\]
\[
\Hom_f^{(2)}(L,M)=\Hom(L,f^{*}M),
\]
\[
\Hom_f^{(3)}(L,M)=\Hom(f_{!}L,M)\quad (\overset{\text{(SGA~4~XVIII~3.1.5)}}{\simeq} \Hom(L,f^{!}M)),
\]
\[
\Hom_f^{(4)}(L,M)=\Hom(f^{!}M,L).
\]
Os elementos de $\Hom_f^{(i)}(L,M)$ serão chamados \emph{flechas de tipo $i$ de $L$ a $M$ sobre $f$} (ou $f$-flechas de tipo $i$ de $L$ a $M$). Os isomorfismos de transitividade dos funtores $f^{*}$, $f_{*}$, $f_{!}$, $f^{!}$ permitem compor as flechas de tipo $i$ para $i$ fixo, e obtêm-se assim certas categorias fibradas ou cofibradas sobre a categoria dos $S$-esquemas: por exemplo, para $i=1$ (resp. $i=3$) obtém-se a categoria fibrada e cofibrada considerada em (SGA~4~XVII~4.1.3) (resp. (SGA~4~XVIII~3.1.13)), cuja fibra em $X$ é a categoria oposta a $D(X)$ (resp. $D(X)$).

\subsection*{1.5}
Sejam, para $i=1,2$, $X_i$ um $S$-esquema, $X=X_1\times_S X_2$, $L_i\in\ob D^{-}(X_i)$. Poremos
\[
L_1\otimes^{L}_{S}L_2=\pr_1^{*}L_1\otimes^{L}\pr_2^{*}L_2,
\]
onde $\pr_i:X\to X_i$ é a projeção canônica.

\subsection*{1.6}
Sejam, para $i=1,2$, $f_i:X_i\to Y_i$ um $S$-morfismo,
\[
f=f_1\times_S f_2:X=X_1\times_S X_2\to Y=Y_1\times_S Y_2,
\]
$L_i\in\ob D^{-}(X_i)$, $M_i\in\ob D^{-}(Y_i)$. O isomorfismo canônico evidente
\begin{equation}\tag{1.6.1}
f_1^{*}M_1\otimes^{L}_{S}f_2^{*}M_2\xrightarrow{\sim}f^{*}(M_1\otimes^{L}_{S}M_2)
\end{equation}
permite definir, para $u_i\in\Hom_{f_i}^{(1)}(L_i,M_i)$ (resp. $\Hom_{f_i}^{(2)}(L_i,M_i)$),
\begin{equation}\tag{1.6.2}
u_1\otimes^{L}_{S}u_2\in\Hom_{f}^{(1)}(L,M)\quad(\text{resp. }\Hom_f^{(2)}(L,M)),
\end{equation}
onde $L=L_1\otimes^{L}_{S}L_2$, $M=M_1\otimes^{L}_{S}M_2$. Se $u_i$, $v_i$ forem componíveis, tem-se
\begin{equation}\tag{1.6.3}
(v_1u_1)\otimes^{L}_{S}(v_2u_2)=(v_1\otimes^{L}_{S}v_2)(u_1\otimes^{L}_{S}u_2).
\end{equation}
Tomando $M_i=f_{i*}L_i$, e $u_i$ de tipo~1, dado pela identidade de $f_{i*}L_i$, 1.6.2 fornece uma \emph{flecha de Künneth} (cf.~SGA~4~XVII~5.4.1.4)
\begin{equation}\tag{1.6.4}
f_{1*}L_1\otimes^{L}_{S}f_{2*}L_2\longrightarrow f_{*}(L_1\otimes^{L}_{S}L_2).
\end{equation}
De acordo com (SGA~$4^{1/2}$ Teorema de finitude 1.9), 1.6.4 é um isomorfismo: de fato, por 1.6.3 nos reduzimos ao caso em que $X_2=Y_2$, $f_2=1$, e se aplica (loc. cit. App.). Quando $f_1$ e $f_2$ são próprios, não é necessário invocar (loc. cit.): 1.6.4 é um caso particular do isomorfismo de Künneth (SGA~4~XVII~5.4.3).

\subsection*{1.7}
Escolhamos uma compactificação $f_i=p_i j_i$, onde $j_i$ é uma imersão aberta e $p_i$ um morfismo próprio. Ao compor o isomorfismo canônico (onde $j=j_1\times_S j_2$)
\[
j_{1!}L_1\otimes^{L}_{S}j_{2!}L_2\xrightarrow{\sim}j_{!}(L_1\otimes^{L}_{S}L_2)
\]
com 1.6.4 para $(p_i,j_{i!}L_i)$, obtém-se um isomorfismo de Künneth (SGA~4~XVII~5.4.3)
\begin{equation}\tag{1.7.1}
f_{1!}L_1\otimes^{L}_{S}f_{2!}L_2\xrightarrow{\sim}f_{!}(L_1\otimes^{L}_{S}L_2).
\end{equation}
Verifica-se facilmente (loc. cit.) que 1.7.1 não depende das compactificações escolhidas e é compatível com os isomorfismos de transitividade para as composições $g_if_i$ (cf.~(loc. cit. 5.2.4)). Graças a 1.7.1, pode-se definir, para $u_i\in\Hom_{f_i}^{(3)}(L_i,M_i)$,
\begin{equation}\tag{1.7.2}
u_1\otimes^{L}_{S}u_2\in\Hom_f^{(3)}(L,M)
\end{equation}
com $L$ e $M$ como em 1.6.2, e tem-se uma fórmula de transitividade análoga a 1.6.3.

Para $L_i=f_i^{!}M_i$, e $u_i$ dado pela identidade de $L_i$, 1.7.2 fornece uma \emph{flecha de Künneth}
\begin{equation}\tag{1.7.3}
f_1^{!}M_1\otimes^{L}_{S}f_2^{!}M_2\longrightarrow f^{!}(M_1\otimes^{L}_{S}M_2).
\end{equation}

\subsection*{Proposição 1.7.4}
A flecha 1.7.3 é um isomorfismo.

Pela fórmula de transitividade, basta tratar o caso em que $X_2=Y_2$, $f_2=1$. Por devissage, reduzimo-nos a $M_2=g_{*}j_{!}A_V$, onde $j:V\to U$ é uma imersão aberta e $g:U\to Y_2$ um morfismo finito, e depois, por (SGA~4~XVIII~3.1.12.3), ao caso em que $M_2=A_{Y_2}$. Por localização sobre $X_2$, pode-se supor, por outro lado, que $f_1=g_1j_1$, onde $j_1$ é uma imersão fechada e $g_1$ um morfismo suave de dimensão relativa pura $n_1$. Por transitividade, basta tratar os casos
\[
\text{a) } f_1=g_1,\qquad \text{b) } f_1=j_1.
\]
No caso a), o morfismo de traço fornece, por adjunção (SGA~4~XVIII~3.2.5), um isomorfismo
\[
f_1^{*}M_1(n_1)[2n_1]\xrightarrow{\sim}f_1^{!}M_1;
\]
esse é compatível com a mudança de base $Y_2\to S$, pois o morfismo de traço também o é (SGA~4~XVIII~2.9), de onde resulta 1.7.4 nesse caso. No caso b), se $i_1$ denota a inclusão do aberto complementar, tem-se uma sequência exata
\[
0\longrightarrow j_{1*}\R j_1^{!}M_1\longrightarrow M_1\longrightarrow \R i_{1*}i_1^{*}M_1\longrightarrow0,
\]
e a compatibilidade de $\R i_{1*}$ com a mudança de base, caso particular de 1.6.4, implica a conclusão.

Graças a 1.7.3, podemos definir, para $u_i\in\Hom_f^{(4)}(L_i,M_i)$,
\begin{equation}\tag{1.7.5}
u_1\otimes^{L}_{S}u_2\in\Hom_f^{(4)}(L,M)
\end{equation}
($L$ e $M$ como em 1.6.2), com uma fórmula de transitividade análoga a 1.6.3.

Por outro lado, com as notações de 1.3, 1.7.3 fornece um isomorfismo
\begin{equation}\tag{1.7.6}
K_{X_1}\otimes^{L}_{S}K_{X_2}\xrightarrow{\sim}K_X.
\end{equation}

\subsection*{1.8}
Sejam $f:X\to Y$ um $S$-morfismo, $L\in\ob D(X)$, $M\in\ob D(Y)$, $S'$ um $S$-esquema, $P\in\ob D(S')$, $f':X'\to Y'$ o morfismo deduzido de $f$ pela mudança de base $S'\to S$. A flecha identidade de $P$ pode ser considerada indistintamente como uma flecha de tipo $i$ ($1\leq i\leq4$) sobre a flecha identidade de $S'$; para $u\in\Hom_f^{(i)}(L,M)$ ($1\leq i\leq4$), pode-se, portanto, segundo 1.6.2, 1.7.2, 1.7.5, considerar
\begin{equation}\tag{1.8.1}
u\otimes^{L}_{S}\Id_P\in\Hom_{f'}^{(i)}(L\otimes^{L}_{S}P,M\otimes^{L}_{S}P).
\end{equation}
Escreveremos às vezes $u\otimes^{L}_{S}P$ em vez de $u\otimes^{L}_{S}\Id_P$.

\section{Funtorialidade de \texorpdfstring{$\uRHom$}{RHom} e produtos tensoriais externos}
